% =====================================================================
\chapter{Principios de conteo y ordenaciones}
\label{tema:1}
% =====================================================================

Contar bien es el primer paso para calcular probabilidades en espacios finitos. Este tema reúne las reglas básicas del recuento (suma, producto, complemento, biyección y división), el principio de inclusión--exclusión para dos y tres conjuntos y el recuento de ordenaciones: variaciones y permutaciones, con y sin repetición. La presentación sigue el capítulo~14, «Cardinality Rules», de Lehman, Leighton y Meyer~\cite{mcs}, texto del curso 6.042J del MIT; el capítulo~1 de Ross~\cite{ross} (§1.2 y §1.3); la sección~1.4, «How to count», de Blitzstein y Hwang~\cite{bh}, y la sección~3.1 de Grinstead y Snell~\cite{gs}.

\section{Reglas de la suma, del producto y del complemento}

Se recuerda que, para un conjunto finito $A$, $\abs{A}$ denota su número de elementos y que $[n]=\{1,2,\dots,n\}$. Contar un conjunto consiste en hallar $\abs{A}$; casi siempre se hace descomponiendo $A$ en piezas más sencillas o poniéndolo en correspondencia con otro conjunto ya contado.

\begin{proposicion}[Regla de la suma]\label{prop:t1-suma}
Si $A_1,\dots,A_m$ son conjuntos finitos disjuntos dos a dos, entonces
\[
\abs{A_1\cup\dots\cup A_m}=\abs{A_1}+\dots+\abs{A_m}.
\]
\end{proposicion}

\begin{proof}
Para $m=2$: si $A_1=\{a_1,\dots,a_p\}$ y $A_2=\{b_1,\dots,b_q\}$ son disjuntos, la lista $a_1,\dots,a_p,b_1,\dots,b_q$ enumera $A_1\cup A_2$ sin repeticiones, luego tiene $p+q$ elementos. El caso general se obtiene por inducción sobre $m$, aplicando el caso $m=2$ a $A_1\cup\dots\cup A_{m-1}$ y $A_m$, que son disjuntos.
\end{proof}

\begin{proposicion}[Regla del producto]\label{prop:t1-producto}
Si $A_1,\dots,A_k$ son conjuntos finitos, $\abs{A_1\times\dots\times A_k}=\abs{A_1}\cdots\abs{A_k}$.
\end{proposicion}

\begin{proof}
Para $k=2$, $A_1\times A_2=\bigcup_{a\in A_1}\bigl(\{a\}\times A_2\bigr)$ es una unión disjunta de $\abs{A_1}$ conjuntos, cada uno con $\abs{A_2}$ elementos; por la regla de la suma, $\abs{A_1\times A_2}=\abs{A_1}\abs{A_2}$. El caso general se demuestra por inducción, identificando $A_1\times\dots\times A_k$ con $(A_1\times\dots\times A_{k-1})\times A_k$.
\end{proof}

En la práctica se usa una versión más flexible, en la que el conjunto de opciones de cada etapa puede depender de las elecciones previas, pero no su número.

\begin{teorema}[Regla del producto generalizada]\label{teo:t1-producto-gen}
Sea $S$ un conjunto de sucesiones $(s_1,\dots,s_k)$ tal que hay $n_1$ valores posibles de $s_1$ y, para cada elección de $s_1,\dots,s_{i-1}$, hay exactamente $n_i$ valores posibles de $s_i$ ($i=2,\dots,k$). Entonces $\abs{S}=n_1n_2\cdots n_k$.
\end{teorema}

\begin{proof}
Por inducción sobre $k$. Si $k=1$, $\abs{S}=n_1$. Supóngase cierto para sucesiones de longitud $k-1$ y sea $S'$ el conjunto de los prefijos $(s_1,\dots,s_{k-1})$ de elementos de $S$; por la hipótesis de inducción, $\abs{S'}=n_1\cdots n_{k-1}$. Agrupando los elementos de $S$ según su prefijo, $S$ es la unión disjunta, sobre los prefijos $p\in S'$, de los conjuntos $\{(p,s_k)\in S\}$, cada uno con $n_k$ elementos. Por la regla de la suma, $\abs{S}=\abs{S'}\,n_k=n_1\cdots n_k$.
\end{proof}

\begin{proposicion}[Regla del complemento]\label{prop:t1-complemento}
Si $A\subseteq U$ y $U$ es finito, $\abs{U\setminus A}=\abs{U}-\abs{A}$.
\end{proposicion}

\begin{proof}
$U$ es la unión disjunta de $A$ y $U\setminus A$; basta aplicar la regla de la suma.
\end{proof}

\begin{ejemplo}[Contraseñas]\label{ej:t1-contrasenas}
Una contraseña está formada por 6 caracteres, cada uno una letra minúscula (26) o un dígito (10). Por la regla del producto hay $36^{6}=2\,176\,782\,336$ contraseñas. Las que contienen \emph{al menos} un dígito se cuentan mejor por complemento: las que no tienen ninguno son $26^{6}=308\,915\,776$, luego hay $36^{6}-26^{6}=1\,867\,866\,560$. Intentar contarlas directamente («elegir la posición del dígito, el dígito y el resto») cuenta varias veces las contraseñas con más de un dígito.
\end{ejemplo}

\section{Reglas de la biyección y de la división}

\begin{proposicion}[Regla de la biyección]\label{prop:t1-biyeccion}
Si existe una biyección $f\colon A\to B$ entre conjuntos finitos, entonces $\abs{A}=\abs{B}$.
\end{proposicion}

La demostración es inmediata: si $A=\{a_1,\dots,a_n\}$ sin repeticiones, entonces $f(a_1),\dots,f(a_n)$ son distintos (inyectividad) y agotan $B$ (sobreyectividad). Esta regla permite contar un conjunto contando otro~\cite[§14.1]{mcs}.

\begin{ejemplo}[Subconjuntos y cadenas de bits]\label{ej:t1-bits}
Sea $A=\{a_1,\dots,a_n\}$. La aplicación que asigna a cada $X\subseteq A$ la cadena $b_1\cdots b_n$ con $b_i=1$ si $a_i\in X$ y $b_i=0$ en otro caso es una biyección entre $\mathcal{P}(A)$ y $\{0,1\}^{n}$. Por las reglas de la biyección y del producto, $A$ tiene $2^{n}$ subconjuntos.
\end{ejemplo}

\begin{proposicion}[Regla de la división]\label{prop:t1-division}
Sea $f\colon A\to B$ una función sobreyectiva entre conjuntos finitos tal que cada $b\in B$ tiene exactamente $d$ preimágenes ($f$ es \emph{$d$ a~1}). Entonces $\abs{A}=d\,\abs{B}$.
\end{proposicion}

\begin{proof}
Los conjuntos $f^{-1}(\{b\})$, con $b\in B$, son disjuntos, cada uno tiene $d$ elementos y su unión es $A$. Por la regla de la suma, $\abs{A}=\sum_{b\in B}d=d\abs{B}$.
\end{proof}

\section{Inclusión--exclusión para dos y tres conjuntos}

Cuando los conjuntos no son disjuntos, la suma de sus cardinales cuenta dos veces los elementos comunes y hay que corregirla.

\begin{teorema}[Inclusión--exclusión]\label{teo:t1-ie}
Para conjuntos finitos $A$, $B$ y $C$:
\begin{align*}
\abs{A\cup B}&=\abs{A}+\abs{B}-\abs{A\cap B},\\
\abs{A\cup B\cup C}&=\abs{A}+\abs{B}+\abs{C}-\abs{A\cap B}-\abs{A\cap C}-\abs{B\cap C}+\abs{A\cap B\cap C}.
\end{align*}
\end{teorema}

\begin{proof}
$A\cup B$ es la unión disjunta de $A$ y $B\setminus A$, y $B$ es la unión disjunta de $A\cap B$ y $B\setminus A$. Por la regla de la suma, $\abs{A\cup B}=\abs{A}+\abs{B\setminus A}=\abs{A}+\abs{B}-\abs{A\cap B}$. Para tres conjuntos se aplica la fórmula anterior a $A\cup B$ y $C$, y se usa la distributiva $(A\cup B)\cap C=(A\cap C)\cup(B\cap C)$:
\begin{align*}
\abs{A\cup B\cup C}&=\abs{A\cup B}+\abs{C}-\abs{(A\cap C)\cup(B\cap C)}\\
&=\abs{A}+\abs{B}-\abs{A\cap B}+\abs{C}-\bigl(\abs{A\cap C}+\abs{B\cap C}-\abs{A\cap B\cap C}\bigr),
\end{align*}
que es la fórmula del enunciado.
\end{proof}

La versión para $m$ conjuntos se demuestra en el \cref{tema:2}, una vez disponibles los números combinatorios.

\begin{ejemplo}\label{ej:t1-divisores}
¿Cuántos enteros de $\{1,\dots,100\}$ son divisibles por 2 o por 3? Sean $A_2$ y $A_3$ los múltiplos de 2 y de 3. Entonces $\abs{A_2}=50$, $\abs{A_3}=\lfloor 100/3\rfloor=33$ y $A_2\cap A_3$ son los múltiplos de 6, con $\lfloor100/6\rfloor=16$ elementos. Hay $50+33-16=67$.
\end{ejemplo}

\section{Variaciones y permutaciones}

Se trata ahora de contar \emph{selecciones ordenadas} de elementos de un conjunto de $n$ elementos. La terminología española distingue según se permita o no repetir.

\begin{definicion}[Variaciones y permutaciones]
Sea $A$ un conjunto con $n$ elementos y $k\ge 0$ un entero.
\begin{enumerate}
  \item Una \emph{variación con repetición} de orden $k$ es una sucesión $(a_1,\dots,a_k)$ de elementos de $A$, posiblemente repetidos; equivale a una función $[k]\to A$.
  \item Una \emph{variación sin repetición} de orden $k\le n$ es una sucesión de $k$ elementos \emph{distintos} de $A$; equivale a una función inyectiva $[k]\to A$.
  \item Una \emph{permutación} de $A$ es una variación sin repetición de orden $n$, es decir, una ordenación de todos sus elementos; equivale a una biyección $[n]\to A$.
\end{enumerate}
Se escribe $n!=1\cdot2\cdots n$, con $0!=1$, y $n^{\underline{k}}=n(n-1)\cdots(n-k+1)$ (\emph{factorial descendente}, con $n^{\underline{0}}=1$).
\end{definicion}

\begin{teorema}\label{teo:t1-variaciones}
Si $\abs{A}=n$, el número de variaciones con repetición de orden $k$ es $n^{k}$; el de variaciones sin repetición de orden $k\le n$ es $n^{\underline{k}}=n!/(n-k)!$, y el de permutaciones es $n!$.
\end{teorema}

\begin{proof}
En el primer caso, cada una de las $k$ posiciones admite $n$ valores y se aplica la regla del producto. En el segundo, $a_1$ admite $n$ valores; elegidos $a_1,\dots,a_{i-1}$, distintos, $a_i$ admite exactamente $n-(i-1)$ valores, sean cuales sean los anteriores. Por la regla del producto generalizada (\cref{teo:t1-producto-gen}), hay $n(n-1)\cdots(n-k+1)$ sucesiones. El tercer caso es el segundo con $k=n$.
\end{proof}

\begin{ejemplo}[Permutaciones circulares]\label{ej:t1-circular}
¿De cuántas formas pueden sentarse $n\ge1$ personas alrededor de una mesa redonda si solo importa quién tiene a quién a su derecha? Cada permutación $(a_1,\dots,a_n)$ determina una disposición circular (sentando a $a_1,a_2,\dots$ en sentido horario a partir de una silla fija), y las $n$ rotaciones de una permutación dan la misma disposición. La aplicación «permutación $\mapsto$ disposición» es, por tanto, $n$ a~1, y por la regla de la división hay $n!/n=(n-1)!$ disposiciones~\cite[§14.4.2]{mcs}.
\end{ejemplo}

\begin{ejemplo}[Funciones entre conjuntos finitos]
Sean $\abs{A}=k$ y $\abs{B}=n$. Hay $n^{k}$ funciones $A\to B$ y $n^{\underline{k}}$ inyectivas. Si $k=n$, las inyectivas son biyectivas y hay $n!$. Por ejemplo, con $k=3$ y $n=5$ hay $125$ funciones, de las que $5\cdot4\cdot3=60$ son inyectivas.
\end{ejemplo}

\begin{observacion}[Conexión con el curso]
Con un vocabulario de $K$ tokens hay $K^{T}$ secuencias de longitud $T$ (variaciones con repetición), y el modelo de lenguaje del módulo~1 asigna una probabilidad a cada una. Con $K=50\,000$ y $T=3$ ya hay $1{,}25\cdot10^{14}$ secuencias, lo que explica por qué una distribución sobre secuencias no puede almacenarse como tabla y se factoriza con la regla de la cadena (\cref{tema:4}). En el módulo~12, el alfabeto de bytes tiene $2^{8}=256$ símbolos y UTF-8 codifica cada carácter con 1 a 4 bytes; el número de cadenas de bytes de cada longitud se cuenta con la regla del producto.
\end{observacion}

\begin{observacion}[Errores frecuentes]
(1) Sumar cuando hay que multiplicar, o al revés: la regla de la suma corresponde a «o» entre casos \emph{disjuntos}; la del producto, a «y» entre etapas sucesivas. (2) Aplicar la regla de la suma a casos que se solapan: hay que usar inclusión--exclusión. (3) Contar «al menos uno» eligiendo primero uno y luego el resto, lo que cuenta varias veces los casos con más de uno; es mejor pasar al complemento. (4) Olvidar si el orden importa: $(a,b)$ y $(b,a)$ son variaciones distintas pero el mismo subconjunto. (5) Aplicar la regla del producto generalizada cuando el número de opciones de una etapa \emph{depende} de las elecciones previas.
\end{observacion}

\begin{lecturas}
\item Lehman, Leighton y Meyer~\cite{mcs}, cap.~14 «Cardinality Rules», §14.1 a §14.4.
\item Ross~\cite{ross}, cap.~1, §1.2 «The Basic Principle of Counting» y §1.3 «Permutations».
\item Blitzstein y Hwang~\cite{bh}, §1.4 «How to count».
\item Grinstead y Snell~\cite{gs}, §3.1 «Permutations».
\item Bertsekas y Tsitsiklis~\cite{bt}, §1.6 (clase~4 de MIT 6.041~\cite{mit6041f10}, «Counting»).
\end{lecturas}

\section*{Ejercicios del tema 1}
\addcontentsline{toc}{section}{Ejercicios del tema 1}

\begin{ejercicio}{1}{eje:1-reglas}
Indique qué regla (suma, producto, complemento, biyección o división) aplicaría en cada caso y dé el resultado: (a) número de menús con un primero entre 4, un segundo entre 5 y un postre entre 3; (b) número de formas de elegir \emph{un} plato, sea primero, segundo o postre; (c) número de cadenas de 10 bits que no son todas ceros; (d) número de subconjuntos de un conjunto de 10 elementos.
\end{ejercicio}

\begin{ejercicio}{1}{eje:1-vocabulario}
Un vocabulario tiene $K=50\,000$ tokens. (a) ¿Cuántas secuencias de 3 tokens hay? (b) ¿Cuántas sin tokens repetidos? (c) ¿Qué fracción de las secuencias de (a) tiene algún token repetido? Dé el resultado con tres cifras significativas.
\end{ejercicio}

\begin{ejercicio}{2}{eje:1-contrasenas}
Una contraseña tiene entre 6 y 8 caracteres, cada uno una letra minúscula (26) o un dígito (10), y debe contener al menos un dígito. ¿Cuántas contraseñas válidas hay?
\fuente{adaptado de Lehman, Leighton y Meyer~\cite{mcs}, §14.2.4 «Counting Passwords».}
\end{ejercicio}

\begin{ejercicio}{2}{eje:1-divisibles}
¿Cuántos enteros de $\{1,\dots,1000\}$ son divisibles por 2, por 3 o por~5? ¿Cuántos no son divisibles por ninguno de los tres?
\end{ejercicio}

\begin{ejercicio}{2}{eje:1-palindromos}
Una \emph{palabra} es cualquier cadena de letras de un alfabeto de 26. (a) ¿Cuántas palabras de 2 letras son palíndromos? ¿Y de 3 letras? (b) Compare la proporción de palíndromos entre las palabras de 2 letras y entre las de 3. (c) Generalice: ¿cuántos palíndromos de longitud $n$ hay?
\fuente{adaptado de Harvard Stat 110~\cite{stat110}, \emph{Strategic Practice}~1, problema~1(b).}
\end{ejercicio}

\begin{ejercicio}{2}{eje:1-mesa}
Ocho personas se sientan alrededor de una mesa redonda; dos disposiciones son iguales si una es rotación de la otra. (a) ¿Cuántas disposiciones hay? (b) ¿En cuántas se sientan juntas dos personas concretas, Ana y Blas? (c) ¿Y en cuántas \emph{no} se sientan juntas?
\fuente{adaptado de Lehman, Leighton y Meyer~\cite{mcs}, §14.4.2 «Knights of the Round Table».}
\end{ejercicio}

\begin{ejercicio}{2}{eje:1-cifras}
¿Cuántos números de cuatro cifras (de 1000 a 9999) tienen todas sus cifras distintas? ¿Cuántos de ellos son pares?
\end{ejercicio}

\begin{ejercicio}{2}{eje:1-funciones}
Sean $A=\{1,2,3,4\}$ y $B=\{x,y\}$. (a) ¿Cuántas funciones $A\to B$ hay? (b) ¿Cuántas son inyectivas? (c) Use el complemento para contar las sobreyectivas. (d) Cuente, mediante inclusión--exclusión, las funciones sobreyectivas de un conjunto de 5 elementos en uno de~3.
\end{ejercicio}

\begin{ejercicio}{2}{eje:1-sombreros}
Se reparten al azar los sombreros de $n$ personas, uno a cada una; un reparto es una permutación de los $n$ sombreros. (a) ¿En cuántos repartos recibe cada persona su sombrero? (b) ¿En cuántos reciben su sombrero las $m$ primeras personas ($m\le n$)? (c) Si $2m\le n$, ¿en cuántos reciben las $m$ primeras personas sombreros que pertenecen a las $m$ últimas?
\fuente{adaptado de MIT 6.041~\cite{mit6041f10}, \emph{Problem Set}~3, problema~1, apartados (a) a~(c).}
\end{ejercicio}

\begin{ejercicio}{2}{eje:1-encuesta}
En un grupo de 60 estudiantes, 30 cursan Cálculo, 25 Álgebra y 20 Programación; 10 cursan Cálculo y Álgebra, 8 Cálculo y Programación, 7 Álgebra y Programación, y 3 las tres. ¿Cuántos no cursan ninguna? ¿Cuántos cursan exactamente una?
\end{ejercicio}

\begin{ejercicio}{3}{eje:1-producto-ind}
Demuestre por inducción sobre $k$ que el número de funciones inyectivas de $[k]$ en un conjunto de $n$ elementos es $n^{\underline{k}}$ para $0\le k\le n$, sin usar la regla del producto generalizada. Deduzca que no hay ninguna si $k>n$ (principio del palomar).
\end{ejercicio}

\begin{ejercicio}{3}{eje:1-fibonacci}
Sea $a_n$ el número de cadenas de $n$ bits sin dos unos consecutivos. Demuestre que $a_1=2$, $a_2=3$ y $a_n=a_{n-1}+a_{n-2}$ para $n\ge3$, justificando cada paso con las reglas de la suma y de la biyección. Calcule $a_8$.
\end{ejercicio}

\begin{ejercicio}{3}{eje:1-ie-cuatro}
A partir del \cref{teo:t1-ie} para tres conjuntos, demuestre la fórmula de inclusión--exclusión para cuatro conjuntos finitos $A,B,C,D$.
\end{ejercicio}

\begin{ejercicio}{3}{eje:1-collares}
(a) Demuestre con la regla de la división que el número de formas de sentar a $n\ge3$ personas alrededor de una mesa, si además se identifican dos disposiciones que son simétricas (una es la otra leída en sentido contrario), es $(n-1)!/2$. (b) Demuestre que el número de formas de repartir $2n$ personas en $n$ parejas (sin orden entre parejas ni dentro de ellas) es $(2n)!/(2^{n}n!)$.
\end{ejercicio}
